\uuid{oXhV}
\exo7id{3467}
\titre{exo7 3467}
\auteur{quercia}
\organisation{exo7}
\datecreate{2010-03-10}
\isIndication{false}
\isCorrection{true}
\chapitre{Déterminant, système linéaire}
\sousChapitre{Système linéaire, rang}
\cpge{mpsi-mp2i.calcul-matriciel-systemes.b, pcsi.calcul-matriciel-systemes-lineaires.b, ptsi.calcul-matriciel-systemes-lineaires.b, mpsi-mp2i.calcul-matriciel-systemes.e, pcsi.determinants.c, ptsi.determinants.a, mpsi-mp2i.determinants.d, pcsi.calcul-matriciel-systemes-lineaires.d, ptsi.calcul-matriciel-systemes-lineaires.d, bcpst1.matrices}
\module{Algèbre}
\niveau{L2}
\difficulte{}

\contenu{
\question{
Soit $A \in \mathcal{M}_n(K)$ inversible et $B$ la matrice obtenue en échangeant dans $A$
les colonnes $i$ et $j$. Montrer que $B$ est aussi inversible. Comment passe-t-on de
$A^{-1}$ à $B^{-1}$ ?
}
\reponse{
Échange des lignes $i$ et $j$.
}
}
